\section{Introducción}

El mercado laboral de Yucatán presenta una dualidad 
estructural característica del sureste mexicano: un 
sector industrial y de servicios en expansión 
coexistiendo con una tasa de informalidad laboral que, 
al cierre del cuarto trimestre de 2024, alcanzó el 
58.1\%\footnote{Dato obtenido de los Indicadores Estratégicos 
de la Encuesta Nacional de Ocupación y Empleo (ENOE), INEGI (2024).}. 
En este contexto, la política de incrementos sustanciales
al salario mínimo real, iniciada en 2019, plantea 
un dilema econométrico relevante.

\subsection{Planteamiento del Problema}
La hipótesis central de este análisis 
sugiere que el incremento del salario mínimo 
no necesariamente genera un despido masivo de 
trabajadores actuales (ajuste de stock), sino 
que altera la probabilidad de que los nuevos 
buscadores de empleo logren insertarse en el 
sector formal (ajuste de flujo). Se postula que:

\begin{equation}
P(\text{Informal}_i) = f(SM_t, Z_i, \epsilon_i)
\end{equation}

Donde la probabilidad de que un individuo $i$ sea informal 
depende del nivel del salario mínimo real ($SM_t$), de un 
vector de características sociodemográficas ($Z_i$) ---como edad 
y escolaridad--- y de un término de error ($\epsilon_i$), el 
cual captura los factores no observados y la naturaleza estocástica 
del comportamiento humano en el mercado laboral.

\subsection{Objetivos}
\begin{itemize}
    \item \textbf{General:} Evaluar el impacto de los incrementos al salario mínimo en la probabilidad de acceso al empleo formal en Yucatán a partir de microdatos de la ENOE.
    \item \textbf{Específico:} Identificar si los jóvenes y trabajadores con menor escolaridad presentan una mayor vulnerabilidad de desplazamiento hacia la informalidad tras los ajustes salariales.
\end{itemize}